\chapter{Comparison Data, Modes, and Interaction Beta Functions}
\label{ch:renormalization}

The main manuscript coarse-grains a finite agent network at \(c_*\) by exact
channels, retaining an effective action, a closure residual, finite-step
multipliers, and a native growth bound
(\cref{ch:coarsegraining,ch:agent-network-rg}). A renormalization flow in the
Wilsonian sense adds maps that compare each blocked theory with the original
units, and only then speaks of modes, exponents, relevance, beta functions,
and fixed points. This part asks which comparison data those words require,
what they license, and which of them a finite network at \(c_*\) can carry.

Comparison data are declared structure independent of the channels, and the
finite networks of the main manuscript do not inhabit the asymptotic
classification: finite levels with different numbers of configurations admit
no stochastic comparison, blocking terminates, and the sign of an exponent
depends on the declared block schedule
(\cref{prop:rg-finite-depth-no-asymptotics}). Where the
classification is available, its canonical content is the forward Lyapunov
filtration (\cref{prop:rg-lyapunov-filtration}); exponents are invariant
only under tempered comparisons
(\cref{thm:rg-tempered-comparison,prop:rg-superexponential-distortion}); and
the exact interaction cocycle has no relevant direction under tempered
comparisons (\cref{cor:rg-interaction-no-relevant-mode}), so a relevant
direction needs an extensive normalization, as in the Gaussian block spectrum
(\cref{def:rg-hermite-relevance}). A beta function is a difference quotient on
one reference space. It depends on the reference trajectory
(\cref{prop:rg-action-beta-reference-change}) and on the comparison scheme
(\cref{prop:rg-scheme-dependence}), and a retained beta differs from the exact
one by the transported closure residual of the main manuscript
(\cref{prop:rg-retained-beta-residual}). These results are used in
\cref{ch:wilsonian-limits}, which adds continuous scale, fixed points,
spectra, and two-index limits.

Throughout this part \(b_k>1\) are declared block ratios, \(\varsigma\) is the
cumulative log block scale, \(J_\ell\) are comparison isomorphisms into the
scale-\(\ell\) objects, \(I_b\) and \(K_b\) are identification and rescaled
kernels, and beta functions are written \(\mathfrak b\), so that they are not
confused with the attention rows \(\beta^x_{ij}\) of the main manuscript. The
measure pairs \((\varrho,\mathsf m)\), the reverse conditional expectations
\(\mathsf U\), the finite-step multipliers \(D_\ell\) and
\(D_{n\leftarrow\ell}\) on Banach spaces \(\mathscr X_\ell\), and the
interaction objects \(\mathcal G_\ell\), \(\mathsf g\), \(\mathcal T_\ell\),
\(\operatorname{asm}_\ell\), \(\operatorname{hd}_\ell\),
\(\mathsf P^{\mathrm{ret}}_\ell\), and \(r^{\mathcal G}_{\ell+1}\) are those of
the main manuscript.

\paragraph{Scope.}
This part asserts no existence theorem for exponents. No Oseledets splitting
of a scale cocycle is claimed: a multiplicative ergodic theorem
\citep{Oseledets1968} requires a measure-preserving base flow, measurability of
the cocycle in the base variable, integrability of \(\log^+\lVert D\rVert\)
and, for the two-sided splitting, of \(\log^+\lVert D^{-1}\rVert\), and either
finite dimension \citep[Chapters~3 and~4]{Arnold1998} or, in infinite
dimensions, compactness of the operators on a Hilbert space
\citep{Ruelle1982} or quasi-compactness on a separable Banach space
\citep{LianLu2010}, and a deterministic scale sequence supplies none of these.
No autonomous spectrum, mode, or exponent is asserted for a nonautonomous
sequence; no continuous beta function is derived from a discrete blocking
sequence; no universality class is identified; and none of the invariant or
fixed objects of this part is identified with a physical law, which would
need the apparatus of \cref{claim:physical-law-identification}.
\status{NOT-CLAIMED}
The open constructions are
\cref{open:rg-rescaling-compatibility,open:rg-multiplicative-ergodic,open:rg-hermite-extensive-relevance,open:rg-hermite-nonlinear-attraction}
in this chapter and
\cref{open:rg-continuous-beta,open:rg-two-index-limits,open:rg-scalarized-attraction}
in \cref{ch:wilsonian-limits}.

\section{Comparison data and dynamical types}
\label{sec:rg-dynamical-types}

\definitionheading{Comparison data and reference cocycle}{def:rg-comparison-data}
Let a scale diagram in the sense of \cref{def:rg-scale-diagram} have objects
\(\mathfrak X_\ell\) and adjacent arrows \(C_{\ell,\ell+1}\) in a declared
target category \(\mathscr K\). Comparison data are a reference object
\(\mathfrak X_\star\) and declared isomorphisms
\(J_\ell:\mathfrak X_\star\xrightarrow{\sim}\mathfrak X_\ell\) in
\(\mathscr K\). Field, base, unit, and reference-measure rescalings are
components of \(J_\ell\), and so is the parameter rescaling of a separately
supplied dynamical model. The reference cocycle is
\begin{equation}
\widehat{\mathfrak R}_\ell
=J_{\ell+1}^{-1}\circ C_{\ell,\ell+1}\circ J_\ell:
\mathfrak X_\star\longrightarrow\mathfrak X_\star,
\label{eq:rg-reference-cocycle}
\end{equation}
where \(\circ\) is composition in \(\mathscr K\); in a stochastic component it
is represented by right-acting kernel juxtaposition
\eqref{eq:rg-right-acting-kernel-composition}. A positive scalar
\(\zeta_\ell\) can be one component of \(J_\ell\) but not all of it, since
multiplication by \(\zeta_\ell\) maps each object to itself and identifies no
two distinct objects.

\definitionheading{Dynamical type}{def:rg-dynamical-type}
The reference cocycle \eqref{eq:rg-reference-cocycle} is autonomous when
\(\widehat{\mathfrak R}_\ell=\widehat{\mathfrak R}\) for every \(\ell\),
periodic of period \(p\) when
\(\widehat{\mathfrak R}_{\ell+p}=\widehat{\mathfrak R}_\ell\), stationary
random when it is a measurable stationary cocycle over a measure-preserving
base, and general otherwise. The invariant object depends on the type. An
autonomous fixed object satisfies \(\widehat{\mathfrak R}x=x\); a periodic
scheme has a fixed object of its monodromy
\(\widehat{\mathfrak R}_{\ell+p-1}\circ\cdots\circ\widehat{\mathfrak R}_\ell\)
or an attracting \(p\)-cycle; a stationary random scheme may have a random
invariant section; a general cocycle may have pullback or forward attracting
sections. A fixed ray additionally requires a positive cone and its
projectivization. An ordered segment
\(\widehat{\mathfrak R}_{n-1}\circ\cdots\circ\widehat{\mathfrak R}_\ell\) is a
two-parameter family, a power of one map only in the autonomous case; writing
it as a power otherwise loses the base scale and, for noncommuting factors,
the order, since
\(A=\begin{psmallmatrix}1&1\\0&1\end{psmallmatrix}\) and
\(B=\begin{psmallmatrix}1&0\\1&1\end{psmallmatrix}\) give
\(BA=\begin{psmallmatrix}1&1\\1&2\end{psmallmatrix}\neq
AB=\begin{psmallmatrix}2&1\\1&1\end{psmallmatrix}\). For a mixed state the
composition law is owed in every declared component category.

\propositionheading{Uniform contraction need not give one fixed ray}{prop:rg-contraction-no-fixed-ray}
If every \(\widehat{\mathfrak R}_\ell\) is a positive map on the interior of a
proper cone that contracts Hilbert's projective metric \(d_H\) by a factor at
most \(\mathrm c<1\), then
\begin{equation}
d_H\bigl(
\widehat{\mathfrak R}_{\ell-1}\cdots\widehat{\mathfrak R}_0x,
\widehat{\mathfrak R}_{\ell-1}\cdots\widehat{\mathfrak R}_0y\bigr)
\leq\mathrm c^\ell d_H(x,y),
\label{eq:rg-cocycle-contraction}
\end{equation}
so initial rays are forgotten, but the attracting section need not be one
scale-independent ray: for the alternating cocycle \(A,B,A,B,\ldots\) on
\(\R^2_{>0}\) with
\(A=\begin{psmallmatrix}2&1\\1&1\end{psmallmatrix}\) and
\(B=\begin{psmallmatrix}1&1\\1&2\end{psmallmatrix}\), even iterates converge
projectively to the ray \([1,\sqrt2]\) and odd iterates to the distinct ray
\(A[1,\sqrt2]\). \status{ESTABLISHED}

\noindent\emph{Proof.}
The bound is the iterated one-step contraction
\citep{Birkhoff1957,nussbaum1988hilbert}. Both matrices have strictly positive
entries, hence finite projective diameter and a common contraction factor
below one. The product
\(BA=\begin{psmallmatrix}3&2\\4&3\end{psmallmatrix}\) has eigenvalues
\(3\pm2\sqrt2\) with Perron eigenvector \((1,\sqrt2)\), so even iterates
\((BA)^kx\) converge projectively to \([1,\sqrt2]\) for every \(x\) in the open
cone, and odd iterates, their images under \(A\), converge to
\((2+\sqrt2,1+\sqrt2)\), which is not proportional to \((1,\sqrt2)\).
\(\square\)

\subsection{The rescaled measure-pair step}

\definitionheading{Rescaled kernel and measure-pair step}{def:rg-rescaled-kernel}
For a blocking ratio \(b>1\), a rescaled step consists of a coarse channel
\(C_b\) and a declared identification kernel \(I_b\) that returns the coarse
space to one common standard Borel space; the rescaled kernel is
\(K_b=C_bI_b\). On measure pairs \((\varrho,\mathsf m)\) in the sense of
\cref{def:rg-measure-pair} the step is
\begin{equation}
\mathfrak R_b(\varrho,\mathsf m)
=\bigl((\varrho C_b)I_b,(\mathsf mC_b)I_b\bigr)
=(\varrho K_b,\mathsf mK_b),
\label{eq:rg-measure-pair-map}
\end{equation}
with reference component \(\mathfrak R^\varrho_b\varrho=\varrho K_b\). The family is
compatible when, after the declared identifications of the ratios,
\begin{equation}
K_{b_1b_2}=K_{b_1}K_{b_2}.
\label{eq:rg-kernel-compatibility}
\end{equation}
Under compatibility
\(\mathfrak R_{b_1b_2}=\mathfrak R_{b_2}\circ\mathfrak R_{b_1}\), because
\((\varrho K_{b_1})K_{b_2}=\varrho K_{b_1b_2}\); otherwise the steps form a cocycle
in the sense of \cref{def:rg-dynamical-type}. Compatibility is the rule
semigroup in which the functor axiom of \cref{def:rg-scale-diagram} has
content.

\remarkheading{A declared family that fails compatibility}{rem:rg-compatibility-computation}
For one declared rescaling family on nested-cycle instances of six categorical
agents, whose identification kernels compose a Wilson-line coarse connection
from the retained dressed boundary generators with a cyclic graph
identification and the identity state identification, blocking by the full
ratio and blocking in declared stages gave coarse laws at total variation
distance \(0.042\) and relative entropy \(0.0050\) (check
\texttt{CHK-RG-RESCALING-COMPATIBILITY}, \cref{tab:base-numerical}). This
family therefore fails \eqref{eq:rg-kernel-compatibility}, and its scale maps
form a cocycle; the computation decides nothing about other families.
\status{NUMERICAL}

\openproblemheading{Compatible rescaling kernels on agent networks}{open:rg-rescaling-compatibility}
The measure-pair steps compose exactly under
\eqref{eq:rg-kernel-compatibility}, the family of
\cref{rem:rg-compatibility-computation} fails that law, and the compatible
cover and path-blocking constructions of
\cref{ch:neighborhood-parents,ch:constructive-spatial-rg} live on their own
domains and supply no identification kernels for the agent network. It is not
known whether identification kernels acting on the gauge-retained network
state of the main manuscript satisfy \eqref{eq:rg-kernel-compatibility} on a
declared class of finite networks at \(c_*\). The problem is settled by such a
family with a proof of \eqref{eq:rg-kernel-compatibility}, or by a proof that
none exists for a stated class, whose composed statements are then stated in
cocycle form.

\section{Generalized cross-scale modes and exponents}
\label{sec:rg-generalized-modes}

A scale sequence supplies no eigenvalue problem. Along an exact orbit it
supplies the finite-step multipliers
\(D_\ell\in\mathcal L(\mathscr X_\ell,\mathscr X_{\ell+1})\) and their ordered products
\(D_{n\leftarrow\ell}\) of \cref{def:rg-derivative-cocycle}, which compose by
\eqref{eq:rg-cocycle-composition} between different Banach spaces
\((\mathscr X_\ell,\lVert\cdot\rVert_\ell)\). This section fixes the mode and exponent
types relative to those declared norms.

\definitionheading{Compatible cross-scale mode line}{def:rg-mode-line}
A rank-one mode line of index \(a\) along the orbit is a family of nonzero
vectors \(v_{\ell,a}\in \mathscr X_\ell\) (in the complexification when \(\mathscr X_\ell\) is
real) and scalars \(\lambda_{\ell,a}\in\mathbb C\) with
\begin{equation}
D_\ell v_{\ell,a}=\lambda_{\ell,a}v_{\ell+1,a}.
\label{eq:rg-mode-line}
\end{equation}
This compatibility equation relates adjacent fibers. The expression
\(D_\ell v=\lambda v\) equates vectors of \(\mathscr X_{\ell+1}\) and \(\mathscr X_\ell\) and is
not used unless an identification is declared: with \(\mathscr X_0=\R\), \(\mathscr X_1=\R^2\),
and \(D_0x=(x,0)\), the line \(v_0=1\), \(v_1=(1,0)\), \(\lambda_0=1\) is well
typed while \(D_0v=\lambda v\) is not. For a mode subbundle of finite rank
\(d_a\), choose injective frames \(V_{\ell,a}:\mathbb C^{d_a}\to \mathscr X_\ell\) and
matrices \(A_{\ell,a}\) with
\begin{equation}
D_\ell V_{\ell,a}=V_{\ell+1,a}A_{\ell,a}.
\label{eq:rg-mode-frame}
\end{equation}

\propositionheading{Ordered product law for mode coefficients}{prop:rg-mode-product}
Under \eqref{eq:rg-mode-line},
\begin{equation}
D_{n\leftarrow\ell}v_{\ell,a}=\lambda_{n\leftarrow\ell,a}v_{n,a},
\qquad
\lambda_{n\leftarrow\ell,a}=\prod_{k=\ell}^{n-1}\lambda_{k,a},
\label{eq:rg-mode-product}
\end{equation}
with the empty product equal to one, and under \eqref{eq:rg-mode-frame} the
finite-step matrix is the ordered product \(A_{n-1,a}\cdots A_{\ell,a}\). A
vanishing factor annihilates the propagated mode from that step on, and the
phase or sign of each factor is retained separately from any growth exponent.
Conversely, the definition constrains nothing by itself: for every
\(u\in \mathscr X_\ell\) whose forward images \(D_{k\leftarrow\ell}u\) are nonzero, the
unit vectors \(v_k=D_{k\leftarrow\ell}u/\lVert D_{k\leftarrow\ell}u\rVert_k\)
and scalars
\(\lambda_k=\lVert D_{k+1\leftarrow\ell}u\rVert_{k+1}/
\lVert D_{k\leftarrow\ell}u\rVert_k\) form a mode line. The index \(a\)
therefore acquires content only from an additional declared constraint, such
as invariance of a declared subbundle, equivariance under a declared symmetry,
or a comparison identification under which the mode section is constant.
\status{ESTABLISHED}

\noindent\emph{Proof.}
Induction on \(n-\ell\) with \eqref{eq:rg-cocycle-composition}:
\(D_{n+1\leftarrow\ell}v_{\ell,a}=D_nD_{n\leftarrow\ell}v_{\ell,a}
=\lambda_{n\leftarrow\ell,a}\lambda_{n,a}v_{n+1,a}\). The frame case is the
same induction with matrix factors in the order fixed by
\eqref{eq:rg-mode-frame}. For the converse,
\(D_kv_k=D_{k+1\leftarrow\ell}u/\lVert D_{k\leftarrow\ell}u\rVert_k
=\lambda_kv_{k+1}\). \(\square\)

The object carried by the multiplier cocycle alone is the forward filtration
of \cref{prop:rg-lyapunov-filtration}.

\definitionheading{Cumulative log block scale, growth rates, and relevance}{def:rg-mode-exponents}
\label{def:rg-action-relevance-conventions}
Let \(b_k>1\) be declared adjacent block ratios and set
\begin{equation}
\Delta\varsigma_k=\log b_k,
\qquad
\varsigma_{n\leftarrow\ell}=\sum_{k=\ell}^{n-1}\Delta\varsigma_k,
\label{eq:rg-cumulative-log-scale}
\end{equation}
with \(\varsigma_{n\leftarrow\ell}\to\infty\). For
\(u\in \mathscr X_\ell\setminus\{0\}\) the finite-time growth rate relative to the
declared norms is
\begin{equation}
\chi_{n\leftarrow\ell}(u)
=\frac{\log\lVert D_{n\leftarrow\ell}u\rVert_n-\log\lVert u\rVert_\ell}
{\varsigma_{n\leftarrow\ell}},
\label{eq:rg-finite-time-rate}
\end{equation}
equal to \(-\infty\) once the image vanishes. The upper rate
\(\bar\chi_\ell(u)\) and the lower rate are its limit superior and limit
inferior as \(n\to\infty\), with \(\bar\chi_\ell(0)=-\infty\); a Lyapunov
exponent is asserted only when they agree. For a mode line the scalar
finite-time exponent is
\begin{equation}
\upsilon_{n\leftarrow\ell,a}
=\frac{\log\lvert\lambda_{n\leftarrow\ell,a}\rvert}{\varsigma_{n\leftarrow\ell}}
=\frac{\sum_{k=\ell}^{n-1}\log\lvert\lambda_{k,a}\rvert}
{\varsigma_{n\leftarrow\ell}},
\qquad\log0=-\infty.
\label{eq:rg-scalar-exponent}
\end{equation}
A nonzero direction \(u\) is relevant, marginal, or irrelevant according as
\(\bar\chi_\ell(u)\) is positive, zero, or negative, and a mode line according
as the limit superior of \(\upsilon_{n\leftarrow\ell,a}\) is positive, zero, or
negative. The classification is relative to the declared norms, the declared
extensive normalization of the mode sections, the declared block ratios, and
the declared comparison maps; it is not a property of the scale sequence
alone. In the autonomous case, where a fixed point, one ratio \(b_k=b\), and
bounded isometric identifications make the linearization one bounded
endomorphism \(T\) of one complex Banach space, an isolated eigenvector
\(T\psi_a=\lambda_a\psi_a\) is a constant mode line with exponent
\(\log\lvert\lambda_a\rvert/\log b\), its sign or phase is retained, and
generalized eigenspaces are used only in finite dimension or for isolated
points of the spectrum with their spectral projections.

\propositionheading{Forward Lyapunov filtration}{prop:rg-lyapunov-filtration}
Under \cref{def:rg-mode-exponents}, set, for \(c\in[-\infty,+\infty]\),
\begin{equation}
V_\ell^{\leq c}=\{u\in \mathscr X_\ell:\bar\chi_\ell(u)\leq c\}.
\label{eq:rg-lyapunov-filtration}
\end{equation}
Then every \(V_\ell^{\leq c}\) is a linear subspace,
\(V_\ell^{\leq c}\subseteq V_\ell^{\leq c'}\) for \(c\leq c'\), and
\begin{equation}
\bar\chi_{\ell+1}(D_\ell u)=\bar\chi_\ell(u)\quad(D_\ell u\neq0),
\qquad
D_\ell V_\ell^{\leq c}\subseteq V_{\ell+1}^{\leq c}.
\label{eq:rg-filtration-invariance}
\end{equation}
If \(\mathscr X_\ell\) is finite dimensional, \(\bar\chi_\ell\) takes at most
\(\dim \mathscr X_\ell\) distinct values on \(\mathscr X_\ell\setminus\{0\}\). The irrelevant
directions are the nonzero vectors of the forward-invariant subspace
\(V_\ell^{<0}=\bigcup_{c<0}V_\ell^{\leq c}\), the marginal ones those of
\(V_\ell^{\leq0}\setminus V_\ell^{<0}\), and relevant directions are canonical
only as nonzero classes of \(\mathscr X_\ell/V_\ell^{\leq0}\); a complementary subspace
of relevant directions is not determined by the forward cocycle.
\status{ESTABLISHED}

\noindent\emph{Proof.}
Since \(\varsigma_{n\leftarrow\ell}\to\infty\), the term
\(\log\lVert u\rVert_\ell/\varsigma_{n\leftarrow\ell}\) vanishes in the limit.
The triangle inequality gives
\(\log\lVert D(u+w)\rVert\leq\log2
+\max\{\log\lVert Du\rVert,\log\lVert Dw\rVert\}\) for
\(D=D_{n\leftarrow\ell}\); after division by \(\varsigma_{n\leftarrow\ell}\)
the constant disappears, and the limit superior of a maximum of two sequences
is the maximum of their limit superiors, so
\(\bar\chi_\ell(u+w)\leq\max\{\bar\chi_\ell(u),\bar\chi_\ell(w)\}\). With
\(\bar\chi_\ell(tu)=\bar\chi_\ell(u)\) for \(t\neq0\), each
\(V_\ell^{\leq c}\) is a subspace, and monotonicity in \(c\) is immediate. For
invariance, \(D_{n\leftarrow\ell+1}D_\ell u=D_{n\leftarrow\ell}u\) and
\(\varsigma_{n\leftarrow\ell}/\varsigma_{n\leftarrow\ell+1}\to1\), so the two
limit superiors agree, including infinite values. If
\(c_1<\cdots<c_k\) are attained at \(u_1,\ldots,u_k\), then
\(u_j\in V_\ell^{\leq c_j}\setminus V_\ell^{\leq c_{j-1}}\), so the subspaces
increase strictly and \(k\leq\dim \mathscr X_\ell\). A nested union of subspaces is a
subspace, and invariance passes to it. For the last assertion take
\(D_k=\operatorname{diag}(2,1/2)\) on \(\R^2\) with \(b_k=e\): the vector
\(e_2\) has rate \(-\log2\) and every vector outside
\(\operatorname{span}(e_2)\) has rate \(\log2\), so every line other than
\(\operatorname{span}(e_2)\) is a complement of relevant directions.
\(\square\)

\propositionheading{Scalar and norm growth agree for tempered mode sections}{prop:rg-mode-tempered-normalization}
For a mode line and every \(n>\ell\) with \(\lambda_{n\leftarrow\ell,a}\neq0\),
\begin{equation}
\chi_{n\leftarrow\ell}(v_{\ell,a})
=\upsilon_{n\leftarrow\ell,a}
+\frac{\log\lVert v_{n,a}\rVert_n-\log\lVert v_{\ell,a}\rVert_\ell}
{\varsigma_{n\leftarrow\ell}}.
\label{eq:rg-mode-exponent-relation}
\end{equation}
The two rates therefore have the same upper and lower limits whenever the mode
sections have unit norm, or more generally whenever
\(\log\lVert v_{n,a}\rVert_n\) is sublinear in \(\varsigma_{n\leftarrow\ell}\).
Under a scalar regauging \(v'_{\ell,a}=c_{\ell,a}v_{\ell,a}\) with
\(c_{\ell,a}\neq0\),
\begin{equation}
\lambda'_{\ell,a}=\lambda_{\ell,a}\frac{c_{\ell,a}}{c_{\ell+1,a}},
\qquad
\lambda'_{n\leftarrow\ell,a}
=\lambda_{n\leftarrow\ell,a}\frac{c_{\ell,a}}{c_{n,a}},
\label{eq:rg-mode-scalar-gauge}
\end{equation}
so the upper and lower limits of the scalar exponent are unchanged whenever
\(\lvert\log\lvert c_{n,a}\rvert\rvert/\varsigma_{n\leftarrow\ell}\to0\). This
tempering condition is sufficient but not necessary for preserving those two
values; when the original scalar exponent converges to a finite limit, it is necessary
for the regauged exponent to converge to that same finite limit. \status{ESTABLISHED}

\noindent\emph{Proof.}
Taking norms in \eqref{eq:rg-mode-product} and substituting into
\eqref{eq:rg-finite-time-rate} gives the remainder as the endpoint norm
difference divided by \(\varsigma_{n\leftarrow\ell}\). Substituting
\(v'_{\ell,a}=c_{\ell,a}v_{\ell,a}\) into \eqref{eq:rg-mode-line} gives the
first equality in \eqref{eq:rg-mode-scalar-gauge}, and telescoping gives the
second, so
\(\upsilon'_{n\leftarrow\ell,a}=\upsilon_{n\leftarrow\ell,a}
+(\log\lvert c_{\ell,a}\rvert-\log\lvert c_{n,a}\rvert)/\varsigma_{n\leftarrow\ell}\),
which proves sufficiency; if \(\upsilon\) and \(\upsilon'\) converge to the
same finite limit, their difference forces
\(\log\lvert c_{n,a}\rvert/\varsigma_{n\leftarrow\ell}\to0\). For an infinite limit
this need not hold: with \(b_k=e\), \(\lambda_{k,a}=e^{-(2k+1)}\), and
\(c_{k,a}=e^k\), the rates at \(\ell=0\) are \(-n\) and \(-n-1\), both tending
to \(-\infty\), while \(\lvert\log c_{n,a}\rvert/n=1\). Necessity also fails
for nonconvergent exponents: take \(\ell=0\), \(b_k=e\), \(c_{0,a}=1\),
\(\log\lvert\lambda_{n\leftarrow0,a}\rvert=nu_n\) and
\(\log\lvert c_{n,a}\rvert=nd_n\), where \((u_n,d_n)\) is \((0,0)\),
\((-1,0)\), \((-1/2,1/2)\), \((-1/2,-1/2)\) for \(n\equiv0,1,2,3\pmod4\). The
regauged exponent \(u_n-d_n\) runs through \(0,-1,-1,0\), with the same limit
superior \(0\) and limit inferior \(-1\) as \(u_n\), while
\(\lvert\log\lvert c_{n,a}\rvert\rvert/n=1/2\) for \(n\equiv2,3\pmod4\).
\(\square\)

For a unit-norm mode line the scalar exponent is exactly the rate of
\(v_{\ell,a}\), so the two clauses of \cref{def:rg-mode-exponents} agree. A
splitting of \(\mathscr X_\ell\) into relevant, marginal, and irrelevant subspaces
requires backward data, such as an exponential dichotomy or the hypotheses of
a multiplicative ergodic theorem (\cref{open:rg-multiplicative-ergodic}).

\theoremheading{Exponent invariance under tempered comparison trivializations}{thm:rg-tempered-comparison}
Let \(J_k:\mathscr X_\star\to \mathscr X_k\) be bounded linear isomorphisms with bounded
inverses, and set
\begin{equation}
\widehat D_k=J_{k+1}^{-1}D_kJ_k,
\qquad
\widehat D_{n\leftarrow\ell}=J_n^{-1}D_{n\leftarrow\ell}J_\ell.
\label{eq:rg-transported-cocycle}
\end{equation}
If the trivializations and their inverses are tempered relative to the
cumulative log block scale,
\begin{equation}
\frac{\log^+\lVert J_n\rVert+\log^+\lVert J_n^{-1}\rVert}{\varsigma_{n\leftarrow\ell}}
\longrightarrow0,
\label{eq:rg-tempered-trivialization}
\end{equation}
then for every \(u\in \mathscr X_\ell\setminus\{0\}\) the upper and lower rates of
\(u\) in the native norms coincide with those of \(J_\ell^{-1}u\) computed for
\(\widehat D_{n\leftarrow\ell}\) in the reference norm. Together with the
scalar tempering condition of \cref{prop:rg-mode-tempered-normalization}, this
preserves both the scalar mode exponent and the norm growth rate.
\status{ESTABLISHED}

\noindent\emph{Proof.}
Write \(x_n=D_{n\leftarrow\ell}u\) and
\(\widehat x_n=J_n^{-1}x_n=\widehat D_{n\leftarrow\ell}J_\ell^{-1}u\), the second
identity in \eqref{eq:rg-transported-cocycle}. Bounded invertibility gives
\begin{equation}
\frac{\lVert x_n\rVert_n}{\lVert J_n\rVert}
\leq\lVert\widehat x_n\rVert_\star
\leq\lVert J_n^{-1}\rVert\lVert x_n\rVert_n.
\label{eq:rg-trivialization-two-sided}
\end{equation}
Taking logarithms changes the numerator of \eqref{eq:rg-finite-time-rate} by
at most \(\log^+\lVert J_n\rVert+\log^+\lVert J_n^{-1}\rVert\), and the fixed
initial scale contributes the constant
\(\lvert\log\lVert J_\ell\rVert\rvert+\lvert\log\lVert J_\ell^{-1}\rVert\rvert\).
Dividing by \(\varsigma_{n\leftarrow\ell}\) and using
\eqref{eq:rg-tempered-trivialization} makes the difference of the two
finite-time rates tend to zero. The scalar statement is
\eqref{eq:rg-mode-scalar-gauge}. \(\square\)

\propositionheading{Superexponential distortion changes the apparent exponent}{prop:rg-superexponential-distortion}
Tempering the scalar mode representatives alone is not sufficient. Let
\(\mathscr X_k=\R\) with the absolute value, every \(D_k\) the identity, \(b_k=e\), and
every scalar gauge factor one, so that the native exponent is zero and the
scalar normalization is tempered. With \(J_ku=e^{k^2}u\), the transported
cocycle is \(\widehat D_{n\leftarrow0}=e^{-n^2}I\) and the apparent reference
rate is \(-n\to-\infty\). The comparison trivialization and its inverse, not
only the mode normalization, must satisfy
\eqref{eq:rg-tempered-trivialization}. \status{ESTABLISHED}

\noindent\emph{Proof.}
\(\widehat D_k=J_{k+1}^{-1}J_k=e^{-2k-1}\), whose product over
\(k=0,\ldots,n-1\) is \(e^{-n^2}\); with \(\varsigma_{n\leftarrow0}=n\),
\eqref{eq:rg-finite-time-rate} gives \(-n\), and
\eqref{eq:rg-tempered-trivialization} fails because
\(\log\lVert J_n\rVert=n^2\). \(\square\)

\openproblemheading{Multiplicative ergodic structure of scale cocycles}{open:rg-multiplicative-ergodic}
For every deterministic multiplier cocycle with
\(\varsigma_{n\leftarrow\ell}\to\infty\), the forward filtration of
\cref{prop:rg-lyapunov-filtration} is canonical, and exponents are invariant
under tempered comparisons (\cref{thm:rg-tempered-comparison}). Existence of
exact exponents and of a splitting into exponent subspaces is not established
for any cocycle of this part, because a deterministic blocking sequence
supplies none of the hypotheses of a multiplicative ergodic theorem listed in
the scope paragraph. The problem is settled by a declared stationary family of
blocking data over a measure-preserving base whose multiplier cocycle meets
those hypotheses, or by a proof that no such family exists for the exact
interaction cocycle of the main manuscript.

\propositionheading{Finite single-point networks do not inhabit the asymptotic classification}{prop:rg-finite-depth-no-asymptotics}
Let \(V_0\) be a finite set of agents at \(c_*\). (i) No stochastic comparison
isomorphism exists between finite levels with different numbers of
configurations. (ii) If every coarse
agent is a block of at least two agents, the depth is at most
\(\lfloor\log_2\lvert V_0\rvert\rfloor\), so the limits \(n\to\infty\) and
the hypothesis \(\varsigma_{n\leftarrow\ell}\to\infty\) of
\cref{def:rg-mode-exponents} are unavailable whatever ratios are declared,
and only the finite-time rates \eqref{eq:rg-finite-time-rate} over the
declared depth are defined; if the spaces \(\mathscr X_\ell\) are finite dimensional
and \(\dim \mathscr X_\ell\) is not constant along the sequence, no linear isomorphisms
\(J_\ell:\mathscr X_\star\to \mathscr X_\ell\) exist, and \cref{thm:rg-tempered-comparison} has
no instance. (iii) The signs of the asymptotic rates depend on the declared
schedule unless it is two-sided bounded: if
\(0<\log b\leq\Delta\varsigma_k,\Delta\varsigma'_k\leq\log B<\infty\), the
upper rates of the two schedules have the same sign and so do the lower
rates, whereas a mode with \(\lambda_k=2\) has exponent \(1\) under
\(b_k=2\) and exponent \(2/(n+1)\to0\) under \(b_k=2^{k+1}\).
\status{ESTABLISHED}

\noindent\emph{Proof.}
(i) Mutually inverse Markov kernels between finite sets are permutation
matrices (\cref{prop:rg-scale-functoriality}). (ii) If
\(\lvert V_{\ell+1}\rvert\leq\lvert V_\ell\rvert/2\) and
\(\lvert V_L\rvert\geq1\), then \(2^L\leq\lvert V_0\rvert\); linear isomorphisms
preserve dimension. For the interaction spaces of the main manuscript the
dimension count and the termination of vertex-ratio blocking are
\cref{prop:rg-interaction-finite-network-obstruction}. (iii) For two bounded
schedules,
\(\varsigma_{n\leftarrow\ell}/\varsigma'_{n\leftarrow\ell}\in
[\log b/\log B,\log B/\log b]\), and the two rates have the same numerator.
Multiplying a sequence by factors in a fixed compact subinterval of
\((0,\infty)\) keeps a positive limit superior positive and a negative one
negative; the relation is symmetric, so a zero limit superior stays zero, and
the same holds for limits inferior. In the example,
\(\varsigma_{n\leftarrow0}=\tfrac12n(n+1)\log2\) while
\(\log\lvert\lambda_{n\leftarrow0}\rvert=n\log2\). \(\square\)

For a finite agent network at \(c_*\), only finite-time rates over the
declared depth are therefore used, and the asymptotic classification of
\cref{def:rg-mode-exponents,prop:rg-lyapunov-filtration,thm:rg-tempered-comparison}
is reserved for the two-index limits of \cref{sec:rg-two-index-limits}, in
which system size grows with depth. An inhabited instance needs unboundedly
many sites and a normalization that identifies the spaces, as in
\cref{def:rg-hermite-relevance}, where \(n\) steps consume \(b^n\) independent
coordinates. Whether a family of finite networks at \(c_*\) has a two-index
limit with nonvacuous comparison maps and nontrivial exponents is part of
\cref{open:rg-two-index-limits}.

\subsection{Contraction certificates and the exact interaction cocycle}

\propositionheading{Dobrushin contraction is a sufficient cocycle certificate}{prop:rg-dobrushin-cocycle}
For a reverse-kernel version of \cref{thm:rg-action-lp-contraction} with
Dobrushin coefficient \(\delta\) of \eqref{eq:rg-action-dobrushin-coefficient},
the reverse conditional expectation induces an operator \(\overline{\mathsf U}\) on the
quotient by constants with, in the oscillation norm
\eqref{eq:rg-action-oscillation-norm},
\begin{equation}
\Vert\overline{\mathsf U}\Vert_{\mathrm{osc}\to\mathrm{osc}}\leq\delta.
\label{eq:rg-action-dobrushin-contraction}
\end{equation}
For a sequence of such arrows with coefficients \(\delta_k\) and declared
ratios \(b_k>1\), write
\begin{equation}
\overline{\mathsf U}_{n\leftarrow\ell}=\overline{\mathsf U}_{n-1}\cdots\overline{\mathsf U}_\ell,
\qquad
\varsigma_{n\leftarrow\ell}=\sum_{k=\ell}^{n-1}\log b_k.
\label{eq:rg-action-cocycle-scale}
\end{equation}
If \(\varsigma_{n\leftarrow\ell}\to\infty\) and
\begin{equation}
\limsup_{n\to\infty}
\frac{\sum_{k=\ell}^{n-1}\log\delta_k}{\varsigma_{n\leftarrow\ell}}<0,
\qquad\log0=-\infty,
\label{eq:rg-action-dobrushin-irrelevance}
\end{equation}
then every nonzero bounded class has negative upper rate in the oscillation
norms, so every direction is irrelevant in the sense of
\cref{def:rg-mode-exponents}; for one ratio \(b\) and \(\delta_k\leq\delta<1\)
the upper rate is at most \(\log\delta/\log b\). The condition is sufficient
only: one-step coefficients equal to one do not prevent complete two-step
forgetting, one-step coefficients below one do not give a negative normalized
rate, and failure of \eqref{eq:rg-action-dobrushin-irrelevance} proves neither
marginality nor relevance. \status{ESTABLISHED}

\noindent\emph{Proof.}
The one-step bound \eqref{eq:rg-action-dobrushin-contraction} and its
iteration are the Dobrushin clauses of \cref{thm:rg-action-lp-contraction}:
\begin{equation}
\Vert\overline{\mathsf U}_{n\leftarrow\ell}\Vert_{\mathrm{osc}\to\mathrm{osc}}
\leq\prod_{k=\ell}^{n-1}\delta_k,
\label{eq:rg-action-dobrushin-product-bound}
\end{equation}
and taking logarithms with the zero convention proves the rate bound. For the
first witness take the reverse matrices
\begin{equation}
R_0=
\begin{pmatrix}
1&0\\0&1\\1/2&1/2
\end{pmatrix},
\qquad
R_1=
\begin{pmatrix}
1/2&1/2&0\\0&0&1
\end{pmatrix},
\label{eq:rg-dobrushin-nonnecessity-witness}
\end{equation}
compatible with the laws \((1/2,1/2)\), \((1/3,1/3,1/3)\), and \((2/3,1/3)\).
Both one-step coefficients are one, but both rows of \(R_1R_0\) are
\((1/2,1/2)\), so the composed coefficient is zero. For the second, take the
law \((1/2,1/2)\) at every scale, let \(J\) have both rows \((1/2,1/2)\), and
set
\begin{equation}
R_k=(1-a_k)I+a_kJ,
\qquad a_k=2^{-k-2}.
\label{eq:rg-dobrushin-one-step-insufficient-witness}
\end{equation}
Each \(R_k\) is a compatible reverse kernel with \(\delta_k=1-a_k<1\), but
\(\prod_k\delta_k>0\), so with fixed \(b_k=b\) the normalized rate is zero.
The identity channel has \(\delta=1\) and leaves every bounded direction
marginal, and a two-cycle has eigenvalue \(-1\), so \(\delta\) carries no
phase information. \(\square\)

\corollaryheading{Temperedness of interaction comparison trivializations}{cor:rg-interaction-tempered}
Suppose declared comparison isomorphisms
\(J_\ell:\mathcal G_\star\to\mathcal G_\ell\) obey
\begin{equation}
\log^+\lVert J_n\rVert+\log^+\lVert J_n^{-1}\rVert
\leq\mathrm c\left(1+\lvert V_n\rvert\right)
\label{eq:rg-interaction-trivialization-growth}
\end{equation}
for a constant \(\mathrm c\) independent of \(n\). This bound holds for every
trivialization assembled from the maps of
\eqref{eq:rg-hoeffding-inverse-and-bounds} together with a bounded gauge
realization, because \(\lVert\operatorname{asm}_\ell\rVert\leq1\) and
\(\lVert\operatorname{hd}_\ell\rVert\leq3^{\lvert V_\ell\rvert}-1\). Then the tempering
hypothesis \eqref{eq:rg-tempered-trivialization} of
\cref{thm:rg-tempered-comparison} holds whenever
\(\lvert V_n\rvert/\varsigma_{n\leftarrow\ell}\to0\), in particular for every
infinite blocking sequence whose vertex counts do not increase and whose
declared ratios satisfy \(b_k\geq b>1\). If those ratios are the vertex ratios
\(\lvert V_k\rvert/\lvert V_{k+1}\rvert\), a finite network admits only
finitely many such steps (\cref{prop:rg-interaction-finite-network-obstruction}),
so no such infinite sequence exists and the clause is vacuous; it has
asymptotic content only for ratios declared independently of the vertex
counts, and every exponent computed with them scales inversely with that
declaration. \status{ESTABLISHED}

\noindent\emph{Proof.}
Insert \eqref{eq:rg-interaction-trivialization-growth} into the left side of
\eqref{eq:rg-tempered-trivialization} and divide by
\(\varsigma_{n\leftarrow\ell}\); the constant term vanishes because
\(\varsigma_{n\leftarrow\ell}\to\infty\). If the vertex counts do not increase,
\(\lvert V_n\rvert\leq\lvert V_\ell\rvert\) while
\(\varsigma_{n\leftarrow\ell}\geq(n-\ell)\log b\to\infty\). The operator bounds
are \eqref{eq:rg-hoeffding-inverse-and-bounds}. \(\square\)

The extraction bound is exponential in \(\lvert V_n\rvert\), so the corollary
gives no tempering when the network grows at least linearly in
\(\varsigma_{n\leftarrow\ell}\), and
\cref{prop:rg-superexponential-distortion} shows that an untempered comparison
can change the reported exponent.

\corollaryheading{No relevant exact interaction mode under tempered comparisons}{cor:rg-interaction-no-relevant-mode}
Assume \(\varsigma_{n\leftarrow\ell}\to\infty\) and
\(\lvert V_n\rvert/\varsigma_{n\leftarrow\ell}\to0\). Along an exact
interaction orbit, with the multipliers \eqref{eq:rg-interaction-mode-generator}
and the native norms of \cref{prop:rg-interaction-native-growth-bound}, every
finite-time rate obeys
\begin{equation}
\chi_{n\leftarrow\ell}(u)\leq
\frac{\lvert V_n\rvert\log3}{\varsigma_{n\leftarrow\ell}},
\label{eq:rg-interaction-rate-bound}
\end{equation}
so every native upper rate is at most zero and no direction is relevant in the
sense of \cref{def:rg-mode-exponents}. If the comparison isomorphisms satisfy
\eqref{eq:rg-tempered-trivialization}, the same holds in the reference norm,
and every mode line of \cref{def:rg-mode-line} for these multipliers whose
sections satisfy the tempering condition of
\cref{prop:rg-mode-tempered-normalization} has scalar exponent with limit
superior at most zero. A positive exponent can arise only from a declared
comparison family or mode normalization that violates these tempering
conditions, that is, from a declared extensive normalization; the exact step
supplies none. \status{ESTABLISHED}

\noindent\emph{Proof.}
If \(D_{n\leftarrow\ell}u=0\) the rate is \(-\infty\). Otherwise
\(\chi_{n\leftarrow\ell}(u)\leq\log\lVert D_{n\leftarrow\ell}\rVert/
\varsigma_{n\leftarrow\ell}\), and \eqref{eq:rg-interaction-native-growth-bound}
with \(\log(3^{\lvert V_n\rvert}-1)<\lvert V_n\rvert\log3\) gives
\eqref{eq:rg-interaction-rate-bound}, whose right side tends to zero. The
reference-norm statement is \cref{thm:rg-tempered-comparison}. For a mode
line, \eqref{eq:rg-mode-exponent-relation} writes the scalar exponent as the
native rate of \(v_{\ell,a}\) minus the endpoint-norm correction, which the
tempering condition sends to zero. \(\square\)

\subsection{An inhabited spectrum: the Gaussian block}

\definitionheading{Relevance conventions for the Gaussian block spectrum}{def:rg-hermite-relevance}
Take the replication sum \eqref{eq:rg-score-replication}, the block statistic
\eqref{eq:rg-gaussian-block-statistic}, the Fisher norm
\(\lVert\cdot\rVert_{L^2(\nu)}\) on \(L^2_0(\nu)\) at every scale, and one
ratio \(b_k=b\), so that \(\Delta\varsigma_k=\log b\) in
\eqref{eq:rg-cumulative-log-scale}. By
\cref{thm:rg-gaussian-hermite-spectrum}, each normalized Hermite function
\(\widehat{\mathrm{He}}_k\) is a constant mode line of the score operator
\eqref{eq:rg-gaussian-score-operator} with coefficient \(b^{1-k/2}\), and its
exponent in the sense of \cref{def:rg-mode-exponents} is
\(\upsilon_k=1-k/2\): degree one is relevant with \(\upsilon_1=\tfrac12\),
degree two is marginal with \(\upsilon_2=0\), and every degree \(k\geq3\) is
irrelevant. The constant mode \(k=0\) would have coefficient \(b\) on
\(L^2(\nu)\) but is absent from the centered tangent space.

The degree-one growth comes from the extensive normalization alone: the score
operator is \(\mathsf U_b\mathscr I_b\) with \(\mathsf U_b\) a contraction
(\cref{thm:rg-action-lp-contraction}) and
\(\lVert\mathscr I_b\rVert=\sqrt b\) by \eqref{eq:rg-score-replication-norm},
while no exact arrow supplies growth
(\cref{prop:rg-interaction-native-growth-bound}).

\propositionheading{Convention dependence and forbidden extensions}{prop:rg-hermite-scope}
The spectrum \eqref{eq:rg-hermite-spectrum} is specific to the scalar
independent Gaussian realization and its normalization, and four boundaries
hold. (i) If the coordinates are jointly centered Gaussian with unit variances
and common pairwise correlation \(\mathrm r>-1/(b-1)\), and the statistic is
\(\bar\xi_{\mathrm r}=(\sum_i\xi_i)/\sqrt{b[1+(b-1)\mathrm r]}\), then
\(\operatorname{Corr}(\xi_i,\bar\xi_{\mathrm r})=\mathrm c_{\mathrm r}=\sqrt{[1+(b-1)\mathrm r]/b}\).
With \(\mathsf U_{b,\mathrm r}\) the reverse conditional expectation under the
correlated block law, the operator
\(\mathscr L_{b,\mathrm r}=\mathsf U_{b,\mathrm r}\mathscr I_b\), applied to the
replicated sum \(\sum_ih(\xi_i)\), which is a centered direction of that law but
not the score of the independent replication of
\cref{prop:rg-score-block-lift}, satisfies
\begin{equation}
\mathscr L_{b,\mathrm r}\widehat{\mathrm{He}}_k=b\mathrm c_{\mathrm r}^k\widehat{\mathrm{He}}_k
=b^{1-k/2}\left[1+(b-1)\mathrm r\right]^{k/2}\widehat{\mathrm{He}}_k,
\label{eq:rg-hermite-correlated}
\end{equation}
which differs from \eqref{eq:rg-hermite-eigenvalues} for every
\(\mathrm r\neq0\); keeping the uncorrelated normalization instead changes the
output law away from \(\nu\), and the operator is then not an endomorphism.
(ii) For independent \(\mathcal N(0,I_d)\) coordinates and componentwise normalized
sums the tensor Hermite functions satisfy
\(\mathscr L_b\widehat{\mathrm{He}}_\alpha=b^{1-\lvert\alpha\rvert/2}\widehat{\mathrm{He}}_\alpha\), so degree \(k\)
has multiplicity \(\binom{d+k-1}{k}\). (iii) A common orthogonal action
commutes with the componentwise normalized sum and preserves each degree
space, whereas a general \(\GL(d)\) transformation changes the reference
covariance and needs an explicitly transported family of laws and tangent
norms; no gauge link, holonomy, or curvature datum enters the scalar theorem.
(iv) \eqref{eq:rg-score-replication} is the score of an independent product
perturbation, a restricted direction in \(L^2_0(\nu^{\otimes b})\); the
theorem says nothing about interaction scores or higher Hoeffding sectors
(\cref{sec:rg-finite-network-interactions}). \status{ESTABLISHED}

\noindent\emph{Proof.}
(i) The pair \((\xi_i,\bar\xi_{\mathrm r})\) is jointly centered Gaussian with unit variances and
correlation \(\mathrm c_{\mathrm r}\), so
\(\xi_i\mid \bar\xi_{\mathrm r}=z\sim \mathcal N(\mathrm c_{\mathrm r} z,1-\mathrm c_{\mathrm r}^2)\); the
generating-function computation \eqref{eq:rg-mehler-generating} with
\(b^{-1/2}\) replaced by \(\mathrm c_{\mathrm r}\) gives
\(\E[\mathrm{He}_k(\xi_i)\mid \bar\xi_{\mathrm r}]=\mathrm c_{\mathrm r}^k\mathrm{He}_k(\bar\xi_{\mathrm r})\), and
summing over \(i\) gives \eqref{eq:rg-hermite-correlated}. (ii) Tensor Hermite
functions factor over coordinates and each factor obeys the scalar identity
with the same \(b^{-1/2}\), so the eigenvalue depends only on the total degree,
and \(\binom{d+k-1}{k}\) multi-indices have total degree \(k\) in \(d\)
variables. (iii) An orthogonal map preserves \(\mathcal N(0,I_d)\) and commutes with the
componentwise sum, while a general invertible \(Q\) sends \(\mathcal N(0,I_d)\) to
\(\mathcal N(0,QQ^\top)\). (iv) is the definition of \(\mathscr I_b\). \(\square\)

\citet[Sections~2, 5 and~7]{JonaLasinio2001} obtains the binary case \(b=2\)
of \eqref{eq:rg-hermite-eigenvalues} by linearizing the induced map on laws at
the centered unit-variance Gaussian, identifies the linearized renormalization
at a self-similar Gaussian fixed point with a conditional expectation, and
gives the two-tangent-space form of the eigenvalue equation with its
composition law across scales, leaving the nonlinear terms of a limit theorem
unpursued. The arbitrary-integer eigenvalues, \cref{lem:rg-dqm-realization},
the Fisher budget \eqref{eq:rg-extensive-fisher-defect}, the exact spectrum
\eqref{eq:rg-hermite-spectrum}, the correlated boundary
\eqref{eq:rg-hermite-correlated}, and the typed cocycle of
\cref{sec:rg-generalized-modes} are not attributed to that source.

\openproblemheading{Extensive relevance beyond the scalar Gaussian block}{open:rg-hermite-extensive-relevance}
\Cref{def:rg-hermite-relevance} inhabits all three relevance classes for
independent scalar Gaussian coordinates under the replication sum, the block
statistic, and a fixed Fisher norm, and \cref{prop:rg-hermite-scope} shows
that correlated coordinates, multivariate covariance, and general \(\GL(d)\)
actions each leave that setting. The corresponding spectrum for a correlated,
anisotropic, or gauge-mixed block is not known. The problem is settled by a
declared transported family of laws and tangent norms for such a block,
together with a diagonalization of its score operator.

\openproblemheading{Nonlinear attraction at the Gaussian block law}{open:rg-hermite-nonlinear-attraction}
The spectrum of \cref{thm:rg-gaussian-hermite-spectrum} controls the
derivative of the normalized block map at \(\nu\). The normalized sum maps
each \(\mathcal N(0,\sigma^2)\) to itself, and the tangent of that curve at
\(\sigma^2=1\) is \(\tfrac12\mathrm{He}_2\), so the marginal degree-two
coefficient is carried by a curve of fixed laws. No remainder estimate for the
nonlinear block map in the score-tangent topology of this part is proved here,
so the spectrum does not by itself give a basin of attraction or the behavior
of the map transverse to that curve. The problem is settled by such an
estimate on a declared neighborhood of \(\nu\).

\section{Beta functions and scheme dependence}
\label{sec:rg-beta-function}

A beta function is a difference quotient formed on one space. It therefore
exists only after the comparison data of \cref{def:rg-comparison-data} or the
rescaled kernel of \cref{def:rg-rescaled-kernel} have been declared.

\definitionheading{Reference-dependent action map and discrete action beta}{def:rg-action-beta}
Fix a rescaled kernel \(K_b\). On a dominated tier the action map keeps its
reference argument,
\begin{equation}
\mathfrak R^H_b[H;\varrho]
:=-\log\frac{d\bigl((e^{-H}\varrho)K_b\bigr)}{d(\varrho K_b)},
\label{eq:rg-reference-dependent-action-map}
\end{equation}
so that
\(\mathfrak R_b(\varrho,e^{-H}\varrho)=(\varrho K_b,e^{-\mathfrak R^H_b[H;\varrho]}\varrho K_b)\).
On a declared vector space of finite-valued actions that is closed under this
map along a declared reference trajectory and carries subtraction, the
discrete action beta is
\begin{equation}
\mathfrak b^H_b[H;\varrho]
=\frac{\mathfrak R^H_b[H;\varrho]-H}{\log b}.
\label{eq:rg-discrete-beta-functional}
\end{equation}
If the evidence mass is deliberately forgotten and finite actions are
quotiented by additive constants, the same formula is read in the quotient.
For extended-valued actions, including points at which the numerator would be
\(+\infty-(+\infty)\), no beta is defined and the measure-pair step
\eqref{eq:rg-measure-pair-map} is the primary object. A beta function on a
truncated coupling vector requires invariance of that finite family; otherwise
the omitted hyperedges contribute an explicit residual
(\cref{prop:rg-retained-beta-residual}).

\propositionheading{Change of reference is an inhomogeneous law for the action beta}{prop:rg-action-beta-reference-change}
Let \(H\) be an action with \(\mathsf m=e^{-H}\varrho\) finite and positive, let
\(\Delta\) be bounded and measurable with \(\varrho(e^{-\Delta})=1\), and set
\begin{equation}
\varrho'=e^{-\Delta}\varrho,
\qquad
H'=H-\Delta,
\qquad\text{so that}\qquad
e^{-H'}\varrho'=e^{-H}\varrho=\mathsf m.
\label{eq:rg-reference-change-pair}
\end{equation}
Then \(\varrho'\) is a probability reference for the same measure, and
\((\varrho'K_b)\)-almost everywhere
\begin{equation}
\mathfrak R^H_b[H';\varrho']
=\mathfrak R^H_b[H;\varrho]-\mathfrak R^H_b[\Delta;\varrho].
\label{eq:rg-action-reference-change}
\end{equation}
On a declared finite-valued action vector space containing \(H\), \(\Delta\),
and their images, the discrete action beta therefore obeys the inhomogeneous
law
\begin{equation}
\mathfrak b^H_b[H';\varrho']
=\mathfrak b^H_b[H;\varrho]-\mathfrak b^H_b[\Delta;\varrho].
\label{eq:rg-beta-reference-change}
\end{equation}
The added term is not a linear reparameterization, and it vanishes exactly
when the increment is itself fixed, \(\mathfrak R^H_b[\Delta;\varrho]=\Delta\). A
beta function is therefore a joint statement about the action and the
declared reference trajectory. \status{ESTABLISHED}

\noindent\emph{Proof.}
Since \(\varrho(e^{-\Delta})=1\) and \(\Delta\) is bounded, \(\varrho'\) is a
probability measure with
\(e^{-\sup\Delta}\varrho\leq\varrho'\leq e^{-\inf\Delta}\varrho\), and
\eqref{eq:rg-reference-change-pair} holds by construction. Pushing these
inequalities through \(K_b\) shows that \(\varrho K_b\) and \(\varrho'K_b\) are
equivalent and that \(d(\varrho'K_b)/d(\varrho K_b)\) lies in
\([e^{-\sup\Delta},e^{-\inf\Delta}]\); by
\eqref{eq:rg-reference-dependent-action-map} it equals
\(\exp(-\mathfrak R^H_b[\Delta;\varrho])\). Both references dominate \(\mathsf mK_b\), and
the chain rule for Radon--Nikodym derivatives gives
\(d(\mathsf mK_b)/d(\varrho'K_b)=\bigl[d(\mathsf mK_b)/d(\varrho K_b)\bigr]
\bigl[d(\varrho'K_b)/d(\varrho K_b)\bigr]^{-1}\) almost everywhere. Taking \(-\log\)
gives \eqref{eq:rg-action-reference-change}; subtracting \(H'=H-\Delta\) and
dividing by \(\log b\) gives \eqref{eq:rg-beta-reference-change}, and the
vanishing criterion is \(\mathfrak R^H_b[\Delta;\varrho]-\Delta=0\). \(\square\)

\definitionheading{Typed discrete interaction beta}{def:rg-interaction-beta}
Let \(\mathcal G_\star\) be a declared reference interaction space and let
\begin{equation}
J_\ell:\mathcal G_\star\xrightarrow{\ \sim\ }\mathcal G_\ell
\label{eq:rg-interaction-trivialization}
\end{equation}
be bounded isomorphisms with bounded inverses, the comparison data of
\cref{def:rg-comparison-data} for the interaction tier. The pulled exact
interaction map and the exact discrete interaction beta are
\begin{equation}
\widehat{\mathcal T}_\ell
=J_{\ell+1}^{-1}\mathcal T_\ell J_\ell:
\mathcal G_\star\longrightarrow\mathcal G_\star,
\qquad
\mathfrak b^{\mathrm{ex}}_\ell(\mathsf g)
=\frac{\widehat{\mathcal T}_\ell(\mathsf g)-\mathsf g}{\Delta\varsigma_\ell},
\qquad
\Delta\varsigma_\ell=\log b_\ell,
\label{eq:rg-interaction-beta}
\end{equation}
with \(\mathcal T_\ell\) the exact step
\eqref{eq:rg-exact-nonlinear-interaction-map}. Both terms of the numerator lie
in \(\mathcal G_\star\); the difference \(\mathcal T_\ell(\mathsf g)-\mathsf g\) of
elements of \(\mathcal G_{\ell+1}\) and \(\mathcal G_\ell\) is not formed. If
the coupling object is a manifold, the difference is replaced by a declared
chart difference, logarithm, or retraction into one tangent space.

The definition has no instance on a finite-alphabet step that reduces the
support cardinality, because the two interaction spaces then have different
dimensions (\cref{prop:rg-interaction-finite-network-obstruction}); on such
steps the untransported residual \eqref{eq:rg-interaction-coordinate-residual}
is used.

\propositionheading{Exact, retained, and residual interaction beta}{prop:rg-retained-beta-residual}
Let \(\mathsf P^{\mathrm{ret}}_\ell\) be the bounded idempotent retained projections
of \eqref{eq:rg-interaction-coordinate-residual}, put
\(\widehat{\mathsf P}_\ell=J_\ell^{-1}\mathsf P^{\mathrm{ret}}_\ell J_\ell\), and define, for
\(\mathsf g\in\operatorname{Ran}\widehat{\mathsf P}_\ell\),
\begin{equation}
\mathfrak b^{\mathrm{ret}}_\ell(\mathsf g)
=\frac{\widehat{\mathsf P}_{\ell+1}\widehat{\mathcal T}_\ell(\mathsf g)-\mathsf g}{\Delta\varsigma_\ell},
\qquad
\delta\mathfrak b_\ell(\mathsf g)
=\mathfrak b^{\mathrm{ex}}_\ell(\mathsf g)-\mathfrak b^{\mathrm{ret}}_\ell(\mathsf g).
\label{eq:rg-retained-beta}
\end{equation}
Then the omitted part of the flow is the transported coordinate residual
divided by the log block step,
\begin{equation}
\delta\mathfrak b_\ell(\mathsf g)
=\frac{(I-\widehat{\mathsf P}_{\ell+1})\widehat{\mathcal T}_\ell(\mathsf g)}{\Delta\varsigma_\ell}
=\frac{J_{\ell+1}^{-1}r^{\mathcal G}_{\ell+1}(J_\ell\mathsf g)}{\Delta\varsigma_\ell},
\label{eq:rg-beta-residual-identity}
\end{equation}
and its action-quotient partner is
\begin{equation}
\delta\overline{\mathfrak b}^{\mathcal A}_{\ell+1}(\mathsf g)
=\frac{\operatorname{asm}_{\ell+1}r^{\mathcal G}_{\ell+1}(J_\ell\mathsf g)}{\Delta\varsigma_\ell},
\qquad
\operatorname{hd}_{\ell+1}\delta\overline{\mathfrak b}^{\mathcal A}_{\ell+1}(\mathsf g)
=\frac{r^{\mathcal G}_{\ell+1}(J_\ell\mathsf g)}{\Delta\varsigma_\ell}.
\label{eq:rg-beta-residual-quotient}
\end{equation}
Consequently \(\delta\mathfrak b_\ell\) vanishes on the retained sector if and
only if the exact image of the retained sector is retained,
\begin{equation}
\delta\mathfrak b_\ell\big|_{\operatorname{Ran}\widehat{\mathsf P}_\ell}=0
\quad\Longleftrightarrow\quad
\mathcal T_\ell(\operatorname{Ran}\mathsf P^{\mathrm{ret}}_\ell)
\subseteq\operatorname{Ran}\mathsf P^{\mathrm{ret}}_{\ell+1},
\label{eq:rg-beta-exactness-criterion}
\end{equation}
which is \eqref{eq:rg-interaction-exact-image-invariance}, and otherwise
\(\delta\mathfrak b_\ell\) obeys the two-sided norm control of
\eqref{eq:rg-interaction-residual-norm-control}. Boundedness and idempotence do
not imply this invariance, and a retained fixed point need not be an exact
one: on \(\mathcal G_\star=\R^2\) with \(b=e\), \(\widehat{\mathsf P}(x,y)=(x,0)\), and
\(\widehat{\mathcal T}(x,y)=(x,x)\), every retained input \((x,0)\) has
\(\mathfrak b^{\mathrm{ret}}=(0,0)\) while
\(\mathfrak b^{\mathrm{ex}}=(0,x)\). \status{ESTABLISHED}

\noindent\emph{Proof.}
Subtracting the two displays of \eqref{eq:rg-retained-beta} cancels \(\mathsf g\) and
leaves \((I-\widehat{\mathsf P}_{\ell+1})\widehat{\mathcal T}_\ell(\mathsf g)/\Delta\varsigma_\ell\).
Substituting the conjugated forms of \(\widehat{\mathsf P}_{\ell+1}\) and
\(\widehat{\mathcal T}_\ell\) gives
\(J_{\ell+1}^{-1}(I-\mathsf P^{\mathrm{ret}}_{\ell+1})\mathcal T_\ell(J_\ell\mathsf g)/
\Delta\varsigma_\ell\), which is \eqref{eq:rg-beta-residual-identity} by the
definition of \(r^{\mathcal G}_{\ell+1}\). Applying \(\operatorname{asm}_{\ell+1}\) and
\(\operatorname{hd}_{\ell+1}\) and using the inverse identities of
\cref{thm:rg-hoeffding-action-isomorphism} with
\eqref{eq:rg-interaction-action-residual} gives
\eqref{eq:rg-beta-residual-quotient}. Since \(J_{\ell+1}^{-1}\) is injective and
\(\Delta\varsigma_\ell>0\), \(\delta\mathfrak b_\ell(\mathsf g)=0\) exactly when
\(r^{\mathcal G}_{\ell+1}(J_\ell\mathsf g)=0\), and \(J_\ell\) maps
\(\operatorname{Ran}\widehat{\mathsf P}_\ell\) onto
\(\operatorname{Ran}\mathsf P^{\mathrm{ret}}_\ell\); quantifying over that sector
and using \eqref{eq:rg-interaction-exact-image-invariance} gives
\eqref{eq:rg-beta-exactness-criterion}. The witness is direct substitution:
\(\widehat{\mathcal T}(x,0)=(x,x)\) and \(\widehat{\mathsf P}(x,x)=(x,0)\). \(\square\)

\subsection{Scheme dependence and invariant candidates}
\label{sec:rg-scheme-dependence}

\propositionheading{Scheme dependence of beta functions}{prop:rg-scheme-dependence}
(i) Let \(g(t)\) be running coordinates with beta function
\(\mathfrak b(g)=dg/dt\). Under a smooth reparameterization \(g'=f(g)\) and a
scale change \(t'=\mathrm ct\) with \(\mathrm c>0\),
\begin{equation}
\mathfrak b'(g')=\frac1{\mathrm c}Df(g)\mathfrak b(g).
\label{eq:rg-beta-change}
\end{equation}
(ii) If the comparison isomorphisms of \cref{def:rg-interaction-beta} are
replaced by \(J'_\ell=J_\ell\mathcal S_\ell\) for bounded isomorphisms
\(\mathcal S_\ell\) of \(\mathcal G_\star\), then
\(\widehat{\mathcal T}'_\ell=\mathcal S_{\ell+1}^{-1}\widehat{\mathcal T}_\ell\mathcal S_\ell\)
and
\begin{equation}
\mathfrak b'_\ell(\mathsf g)
=\frac{\mathcal S_{\ell+1}^{-1}\widehat{\mathcal T}_\ell(\mathcal S_\ell\mathsf g)-\mathsf g}
{\Delta\varsigma_\ell},
\label{eq:rg-interaction-beta-scheme-change}
\end{equation}
which reduces to the pushforward relation
\(\mathfrak b'_\ell(\mathsf g)=\mathcal S^{-1}\mathfrak b^{\mathrm{ex}}_\ell(\mathcal S\mathsf g)\)
only for a scale-independent \(\mathcal S_\ell\equiv\mathcal S\); in a moving
comparison frame even the identity step has nonzero beta: on
\(\mathcal G_\ell=\R\) with \(\mathcal T_\ell=\operatorname{id}\), \(b_\ell=e\), and
\(J_\ell u=a_\ell u\), \eqref{eq:rg-interaction-beta} gives
\(\mathfrak b_\ell(\mathsf g)=(a_\ell/a_{\ell+1}-1)\mathsf g\). (iii) For a differentiable
positive endomorphism homogeneous of degree one with a fixed ray, radial
eigenvalue \(\mu_0\), and transverse eigenvalue \(\mu_a\), define for a
declared block factor \(b\)
\begin{equation}
r_a=\frac{\mu_a}{\mu_0},\qquad
\upsilon_a=\frac{\log\lvert r_a\rvert}{\log b}.
\label{eq:rg-projective-dimensions}
\end{equation}
The ratios \(r_a\) are the eigenvalues at the fixed ray of the derivative of
the induced map on rays, and \(r_a\) and \(\upsilon_a\) are invariant under
conjugation of the induced map by a \(C^1\) diffeomorphism near the fixed ray
with invertible derivative there, provided the same \(b\) is used and, when a
declared gauge group has an orbit of fixed rays through the ray, the conjugacy
maps that orbit to the corresponding orbit and transverse eigenvalues are
computed on the quotient by its tangent space. Invertibility of the derivative
is needed: \(x\mapsto x^3\) conjugates \(x\mapsto2x\) to \(y\mapsto8y\).
\status{ESTABLISHED}

\noindent\emph{Proof.}
(i) By the chain rule,
\(dg'/dt'=Df(g)(dg/dt)(dt/dt')=\mathrm c^{-1}Df(g)\mathfrak b(g)\). (ii)
Substitution gives
\(\widehat{\mathcal T}'_\ell=\mathcal S_{\ell+1}^{-1}J_{\ell+1}^{-1}\mathcal T_\ell
J_\ell\mathcal S_\ell\) and then \eqref{eq:rg-interaction-beta-scheme-change};
for constant \(\mathcal S\) the numerator is
\(\mathcal S^{-1}(\widehat{\mathcal T}_\ell(\mathcal S\mathsf g)-\mathcal S\mathsf g)\). In
the example \(\widehat{\mathcal T}_\ell=J_{\ell+1}^{-1}J_\ell=a_\ell/a_{\ell+1}\) and
\(\Delta\varsigma_\ell=1\). (iii) Differentiating \(F(tx)=tF(x)\) in \(t\)
gives Euler's relation \(DF(x)x=F(x)\), so at a fixed ray \(F(x_*)=\mu_0x_*\)
the ray is an eigendirection of \(DF(x_*)\) with eigenvalue \(\mu_0\), and in
an affine chart of rays the derivative of the induced map is \(DF(x_*)/\mu_0\)
on the quotient by the ray, with eigenvalues \(r_a\). Under a \(C^1\)
conjugacy with invertible derivative at the ray, the chain rule makes the
conjugated derivative similar to the original; a conjugacy mapping the orbit
of fixed rays to the corresponding orbit carries the tangent space of the
orbit to its counterpart, so the induced quotient maps are similar as well.
For the last claim, \((2y^{1/3})^3=8y\), and the derivative of \(x^3\) vanishes
at the fixed point. \(\square\)

Raw couplings, beta components, and a chosen blocking statistic are
therefore scheme-dependent coordinates; they are diagnostics once their scheme
is stated. Candidate cross-scheme quantities are conjugacy classes, invariant
manifolds, projective ratios, dimensionless amplitude ratios, and matched
spectral or heat exponents. Their existence, convergence, and completeness in a
given realization, on which any universality statement rests, are part of
\cref{open:rg-two-index-limits}.
